\begin{frame}
  \frametitle{Outline}
\nocite{Bou18tra}
\nocite{CJP15alg}
\nocite{CST10sto}
\nocite{CKS14pri}
\nocite{GPFA13lim}
Most instruments trade on \emph{order-driven} markets:
the participants supply the prices,
and the venue only matches them.

We describe the object that produces the prices we later model:
a pair of queues of resting orders,
sorted by a rule,
matched by an algorithm.

We then look at how that object changes over time,
and at the cost this imposes on anyone who trades size.
\end{frame}

\subsection{Order-driven markets}

\begin{frame}
  \frametitle{Two kinds of order}
A \emph{limit order} carries a price.
It transacts at that price or better, never worse,
and it waits if no counterpart is available.
It \emph{supplies} liquidity.

\pause

A \emph{market order} carries no price.
It transacts immediately, at whatever price the book offers.
It \emph{consumes} liquidity.

\pause

The sender of a limit order chooses the price but not the time of execution;
the sender of a market order chooses the time but not the price.
Nearly everything that follows is a consequence of this asymmetry.
\end{frame}

\subsection{The limit order book}

\begin{frame}
  \frametitle{The book}
\begin{defi}
\label{def.limitOrderBook}
The \emph{limit order book} at time $t$ is the pair of functions
\[
p \mapsto \bidVolumePriceP, \qquad p \mapsto \askVolumePriceP, \qquad p \in \tickSizeOfLOB \Z,
\]
giving the total size of the buy and of the sell limit orders resting at price $p$.
\end{defi}

\pause

Two queues facing one another across a gap:
the bid side at the low prices,
the ask side at the high prices.
Prices live on the grid set by the tick size $\tickSizeOfLOB$.
\end{frame}

\begin{frame}
  \frametitle{Top of the book}
\begin{defi}
\label{def.bestQuotes}
The \emph{best bid} and \emph{best ask} prices at time $t$ are
\[
\bestBidPrice_t := \max \lbrace p : \bidVolumePriceP > 0 \rbrace,
\qquad
\bestAskPrice_t := \min \lbrace p : \askVolumePriceP > 0 \rbrace,
\]
and $\bestBidVolume_t$, $\bestAskVolume_t$ are the volumes resting there.
\end{defi}

\pause

The sides never overlap, so the \emph{bid-ask spread} satisfies
\[
\LOBspread_t = \bestAskPrice_t - \bestBidPrice_t \geq \tickSizeOfLOB.
\]
It is the price of immediacy:
buy at the ask and sell at the bid, and you have paid $\LOBspread_t$ per unit
for not having waited.
\end{frame}

\begin{frame}
  \frametitle{Mid-price and imbalance}
The \emph{mid-price}
\[
\midPrice_t = \frac{\bestBidPrice_t + \bestAskPrice_t}{2}
\]
is the usual answer to ``what is the price?'',
though nobody can trade at it.

\pause

The \emph{queue imbalance}
\[
\volumeImbalance_t = \frac{\bestBidVolume_t - \bestAskVolume_t}{\bestBidVolume_t + \bestAskVolume_t} \in [-1,1]
\]
measures how lopsided the top of the book is.
It is among the most robust short-horizon predictors of the next mid-price move.
\end{frame}

\begin{frame}
  \frametitle{Execution priority}
\begin{defi}
\label{def.executionPriority}
A resting buy limit order with price $p$ submitted at time $s$
has \emph{higher priority} than one with price $p\derivative$ submitted at $s\derivative$ if
\[
p > p\derivative,
\qquad \text{ or } \qquad
p = p\derivative \, \text{ and } \, s < s\derivative.
\]
\end{defi}

\pause

Price-time priority makes a limit order a queueing problem.
Your position in the queue is an asset,
and cancelling forfeits it.
\end{frame}

\subsection{Order flow}

\begin{frame}
  \frametitle{Three events, and nothing else}
A book changes only by
\emph{i.} submission of a new limit order;
\emph{ii.} cancellation of a resting one;
\emph{iii.} execution against an incoming market order.

\pause

Per price level, with $\arrivals$ counting volume added and $\departures$ volume removed,
\[
\bidQueue = \bidQueue[0] + \bidArrivals - \bidDepartures.
\]
An accounting identity; it becomes a model once the law of $\arrivals$ and $\departures$
is specified.

\pause

Poisson arrivals give a tractable model \citep{CST10sto},
but real order flow is clustered,
which calls for self- and cross-exciting intensities $\intensity$.
\end{frame}

\begin{frame}
  \frametitle{Order flow imbalance}
Aggregate the signed changes at the best quotes over $(s,t]$:
\[
\OFI_{s,t} =
\big(\bidArrivals[t] - \bidArrivals[s]\big)
- \big(\bidDepartures[t] - \bidDepartures[s]\big)
- \big(\askArrivals[t] - \askArrivals[s]\big)
+ \big(\askDepartures[t] - \askDepartures[s]\big).
\]

\pause

Contemporaneous mid-price changes are close to linear in $\OFI_{s,t}$,
with slope scaling as the inverse of the depth \citep{CKS14pri}.

This is one of the cleanest empirical regularities in microstructure,
and a good first exercise on real data.
\end{frame}

\begin{frame}
  \frametitle{Price impact}
A market order to buy consumes the ask queue upwards:
each unit is bought worse than the last,
and the best ask ends up higher.

\pause

Part of the move persists --- others read the trade as information ---
and part decays as the level is replenished.

\pause

Hence a large order is split into many small ones released over time.
The trade-off is the one we started from:
trading slowly reduces impact,
but increases exposure to price risk over the execution window.
\end{frame}

\begin{frame}
  \frametitle{What to take away}
There is no such thing as \emph{the} price:
best bid, best ask, mid, last traded --- they differ,
and the wrong choice quietly corrupts a calculation.

\pause

Liquidity is finite and is consumed,
so the cost of a trade is a function of its size.

\pause

The book is a queueing system.
Counting processes and their intensities are the natural tools,
which is where this strand of the course meets the rest of it.
\end{frame}
